\section{Basis vectors as partial derivatives and Jacobian}

In two-dimensional Euclidean space, a point is described by Cartesian coordinates
\[
(x,y)
\]

The position vector is
\[
\mathbf{R} = x \begingroup\color{blue}\vec{e_x}\endgroup + y \begingroup\color{blue}\vec{e_y}\endgroup
\]
where the basis vectors are
\[
\begingroup\color{blue}\vec{e_x}\endgroup = \begin{pmatrix}1 \\ 0\end{pmatrix},
\qquad
\begingroup\color{blue}\vec{e_y}\endgroup = \begin{pmatrix}0 \\ 1\end{pmatrix}.
\]

These satisfy:
\[
\begingroup\color{blue}\vec{e_x}\endgroup \cdot \begingroup\color{blue}\vec{e_x}\endgroup = 1, 
\quad
\begingroup\color{blue}\vec{e_y}\endgroup \cdot \begingroup\color{blue}\vec{e_y}\endgroup = 1,
\quad
\begingroup\color{blue}\vec{e_x}\endgroup \cdot \begingroup\color{blue}\vec{e_y}\endgroup = 0.
\]

Thus, $\{\begingroup\color{blue}\vec{e_x}\endgroup, \begingroup\color{blue}\vec{e_y}\endgroup\}$ is an orthonormal basis.

A key feature:

\[
\frac{\partial \begingroup\color{blue}\vec{e_x}\endgroup}{\partial x} = 
\frac{\partial \begingroup\color{blue}\vec{e_x}\endgroup}{\partial y} =
\frac{\partial \begingroup\color{blue}\vec{e_y}\endgroup}{\partial x} =
\frac{\partial \begingroup\color{blue}\vec{e_y}\endgroup}{\partial y} = 0.
\]

The basis vectors are constant everywhere in space.\\

We now define polar coordinates $(r,\theta)$ by the transformation

\[
x = r \cos\theta,
\qquad
y = r \sin\theta.
\]

The position vector becomes

\[
\mathbf{R} = r\cos\theta \,\begingroup\color{blue}\vec{e_x}\endgroup + r\sin\theta \,\begingroup\color{blue}\vec{e_y}\endgroup
\]

We define the radial unit vector as

\[
\begingroup\color{red}\tilde{\vec{e_r}}\endgroup
= \cos\theta \,\begingroup\color{blue}\vec{e_x}\endgroup 
+ \sin\theta \,\begingroup\color{blue}\vec{e_y}\endgroup.
\]

This points directly away from the origin.\\
And then define the angular unit vector as

\[
\begingroup\color{red}\tilde{\vec{e_\theta}}\endgroup
= -\sin\theta \,\begingroup\color{blue}\vec{e_x}\endgroup
+ \cos\theta \,\begingroup\color{blue}\vec{e_y}\endgroup
\]

This points in the direction of increasing $\theta$.\\
These vectors satisfy

\[
\begingroup\color{red}\tilde{\vec{e_r}}\endgroup \cdot \begingroup\color{red}\tilde{\vec{e_r}}\endgroup = 1,
\quad
\begingroup\color{red}\tilde{\vec{e_\theta}}\endgroup \cdot \begingroup\color{red}\tilde{\vec{e_\theta}}\endgroup = 1,
\quad
\begingroup\color{red}\tilde{\vec{e_r}}\endgroup \cdot \begingroup\color{red}\tilde{\vec{e_\theta}}\endgroup = 0
\]

So they form an orthonormal basis at each point, but polar basis vectors depend on $\theta$\\

Compute derivatives:

\[
\frac{\partial \begingroup\color{red}\tilde{\vec{e_r}}\endgroup}{\partial \theta}
=
-\sin\theta\,\begingroup\color{blue}\vec{e_x}\endgroup
+
\cos\theta\,\begingroup\color{blue}\vec{e_y}\endgroup
=
\begingroup\color{red}\tilde{\vec{e_\theta}}\endgroup
\]

\[
\frac{\partial \begingroup\color{red}\tilde{\vec{e_\theta}}\endgroup}{\partial \theta}
=
-\cos\theta\,\begingroup\color{blue}\vec{e_x}\endgroup
-
\sin\theta\,\begingroup\color{blue}\vec{e_y}\endgroup
=
-\begingroup\color{red}\tilde{\vec{e_r}}\endgroup
\]

Thus,

\[
\frac{\partial \begingroup\color{red}\tilde{\vec{e_r}}\endgroup}{\partial r} = 0,
\qquad
\frac{\partial \begingroup\color{red}\tilde{\vec{e_\theta}}\endgroup}{\partial r} = 0
\]

Key result:

\[
\frac{\partial \begingroup\color{red}\tilde{\vec{e_r}}\endgroup}{\partial \theta} = \begingroup\color{red}\tilde{\vec{e_\theta}}\endgroup,
\qquad
\frac{\partial \begingroup\color{red}\tilde{\vec{e_\theta}}\endgroup}{\partial \theta} = -\begingroup\color{red}\tilde{\vec{e_r}}\endgroup
\]

The basis vectors rotate as $\theta$ changes.\\

Now if you think about this, to find out the directions of the basis vectors, what we do is
following the directions of the coordinate lines and in the case of polar coordinates, for the $\begingroup\color{red}\tilde{\vec{e_\theta}}\endgroup$
basis vector, since the coordinate lines are actually curved, we need to pick the straightest possible
direction when moving along it, which results in being the vector tangent to the curve.\\

\textbf{This reasoning gives us a hint for the next important reinterpretation of basis vectors, as represented by a 
partial derivative of the position vector R}

\begin{equation}
    \begingroup\color{blue}\vec{e_x}\endgroup \equiv \frac{\partial \mathbf{R}}{\partial \begingroup\color{blue}x\endgroup}
\end{equation}

If you carefully think about it, it might look weird at first, but it will eventually make sense.
Think about $\mathbf{R}$ as a function of the coordinates x and y for example, and then
consider adding a small value $\Delta x$ to the x coordinate. The difference between these two position
vectors results in a vector that is parallel to the x coordinate axis.
So this resultant vector has the direction of the basis vector $\color{blue}\vec{e_x}$. \\
This observation is interesting, but we need to solve tha fact that the magnitude of this vector
is varying , depending on the value of the $\Delta x$ used. \\
This can be solved by dividing the resultant vector by the quantity $\Delta x$.
By doing so, we obtain a vector pointing in the x-basis vector direction and magnitude 1...now we 
can proceed further by taking the limit as $\Delta x$ approaches to 0:

\begin{equation}
    \frac{\partial \mathbf{R}}{\partial x} = \lim_{\Delta x \to 0}\frac{\mathbf{R(x + \Delta x, y)} - \mathbf{R(x,y)}}{\Delta x}
\end{equation}

Since this derivative always has a length of one and points the x-axis, we can regard it as the basis vector:

\begin{equation}
    \begingroup\color{blue}\vec{e_x}\endgroup \equiv \frac{\partial \mathbf{R}}{\partial \begingroup\color{blue}x\endgroup}
\end{equation}

This reasoning can be easily applied to the other direction, but generalized further
to any coordinate system.\\
The neat thing about this construction is that you can do it for any position vector, and you 
will still get your basis vectors. Fundamentally, this trick work because, in obtaining a slightly 
off position vector, a small change in coordinate is added in one coordinate only while keeping other coordinates constant. 
Then, we divide the resultant vector (obtained from the difference between the original position vector and the new slightly off one) 
by that small change to normalize the vector.\\\\

Now you should remember when we discussed, in the other notes for "Tensors for Beginners" about the 
forward and backward transformations when we move from one basis vector coordinate system to another.\\
Turns out that those transformation rules can be obtained easily once we introduced this above
concept about basis vectors being partial derivatives of the position vector with respect to the 
coordinate variables.\\
And also, that these transformation matrices have a proper name: \textbf{the Jacobian matrices}\\

Let's say we want to find the forward transformation matrix (the Jacobian) matric when moving 
from cartesian to polar coordinates. This would be just a matter of applying the chain rule:

\begin{align*}
    \begingroup\color{red}\tilde{\vec{e_r}}\endgroup &= \frac{\partial R}{\partial r} = \frac{\partial R}{\partial x}\frac{\partial x}{\partial r} + \frac{\partial R}{\partial y}\frac{\partial y}{\partial r} = \frac{\partial x}{\partial r}\begingroup\color{blue}\vec{e_x}\endgroup + \frac{\partial y}{\partial r}\begingroup\color{blue}\vec{e_y}\endgroup\\
    \begingroup\color{red}\tilde{\vec{e_\theta}}\endgroup &= \frac{\partial R}{\partial \theta} = \frac{\partial R}{\partial x}\frac{\partial x}{\partial \theta} + \frac{\partial R}{\partial y}\frac{\partial y}{\partial \theta} = \frac{\partial x}{\partial \theta}\begingroup\color{blue}\vec{e_x}\endgroup + \frac{\partial y}{\partial \theta}\begingroup\color{blue}\vec{e_y}\endgroup
\end{align*}

We know the transformation rules between cartesian and polar coordinates:

\[
x = r \cos\theta,
\qquad
y = r \sin\theta.
\]

so it's a matter of taking the derivatives to find that the Jacobian matrix, or forward matrix is:

\begin{align*}
    \frac{\partial R}{\partial r} &= \frac{\partial R}{\partial x}\frac{\partial x}{\partial r} + \frac{\partial R}{\partial y}\frac{\partial y}{\partial r} = cos\theta \frac{\partial R}{\partial x} + sin\theta \frac{\partial R}{\partial y}\\
    \frac{\partial R}{\partial \theta} &= \frac{\partial R}{\partial x}\frac{\partial x}{\partial \theta} + \frac{\partial R}{\partial y}\frac{\partial y}{\partial \theta} = -r\sin\theta \frac{\partial R}{\partial x} + r\cos\theta \frac{\partial R}{\partial y}
\end{align*}

So:

\begin{equation}
    F =
    \begin{bmatrix}
        \frac{\partial x}{\partial r} & \frac{\partial x}{\partial \theta}\\
        \frac{\partial y}{\partial r} & \frac{\partial y}{\partial \theta}
    \end{bmatrix} = 
    \begin{bmatrix}
        cos\theta & -r\sin\theta\\
        \sin\theta & r\cos\theta
    \end{bmatrix}
\end{equation}

And you can tackle, with the same derivation logic, the inverse Jacobian matrix, of backward transformation 
matrix, to obviously find out that the multiplication of the two yields the identity matrix.\\

Now, following a more tensorial approach, we can also re-write these transformation rules as:

\begin{align*}
    \frac{\partial \begingroup\color{violet}R\endgroup}{\partial \begingroup\color{magenta}p^i\endgroup} &= \frac{\partial \begingroup\color{violet}R\endgroup}{\partial \begingroup\color{cyan}c^j\endgroup}\frac{\partial \begingroup\color{cyan}c^j\endgroup}{\partial \begingroup\color{magenta}p^i\endgroup} \\
    \frac{\partial \begingroup\color{violet}R\endgroup}{\partial \begingroup\color{cyan}c^i\endgroup} &= \frac{\partial \begingroup\color{violet}R\endgroup}{\partial \begingroup\color{magenta}p^j\endgroup}\frac{\partial \begingroup\color{magenta}p^j\endgroup}{\partial \begingroup\color{cyan}c^i\endgroup} 
\end{align*}


